\section{Build the arithmetic beneath a modular power}
RSA arithmetic repeatedly raises an integer to a power and takes a remainder. For
example, \(65^{17}\bmod3233=2790\). Computing the entire power first would produce a
large intermediate value even in this small example. Reducing after each multiplication
keeps the stored values below the modulus, but the multiplication itself can still
need twice as many bits as either input. This chapter follows that width problem from
one register to a value spread across four words.

Begin with a greatest common divisor and a small modular power. Then represent an
unsigned 128-bit integer as four 32-bit pieces. Write a full product to understand
how partial products and carries fit together. Next build addition, reduction and
multiplication modulo a 128-bit number. Use those operations to raise a 128-bit base
to a full 128-bit exponent. Compare the results with an independent arbitrary-precision
calculation rather than with another transcription of the assembly.

The size is a teaching boundary. Four limbs make every carry and byte position
observable in the debugger; they do not provide a secure RSA key size. The finished
arithmetic functions are also not a complete encryption or signature implementation.
After checking their results, we will identify the additional encoding and scheme
requirements that separate a modular power from those operations.

At each boundary, keep a small failing calculation. If a product is wrong in limb
three, inspect the carries that reach limb three before debugging exponentiation.
If products and remainders agree but powers fail, inspect which exponent bit selected
the next multiply and which register or stack slot retained the current result.
